\clearpage
\section[Conclusão]{Conclusão}

Como minha segunda experiência de projeto na Ages, estou totalmente
confiante para dizer que o ambiente em que estava inserido era um dos melhores de
desenvolvimento que já estive. Acredito que, apesar da complexidade do projeto e
da dificuldade da equipe de lidar com algo desconhecido, que gastou muito recurso
humano, a entrega da equipe foi um produto de valor e que serviu muito bem como
ferramenta de aprendizado.

Achei muito interessante trabalhar com cada um dos meus colegas de equipe
e até sai de lá com várias amizades, tendo em vista o ótimo ambiente de trabalho
criado tanto pelo professor orientador, quanto pelos colegas. Acredito que houve
muitas trocas de informações valiosas e muitos momentos do qual me lembrarei
para a minha vida nessa Ages.

Além disso, fiquei muito de feliz de conseguir integrar conhecimentos de
matérias que estou fazendo em paralelo, para resolver tarefas da Ages, tendo como
um exemplo, a matéria de Verificação e Validação de Software, que ensinou a
realizar testes unitários, de integração e de sistema, que foram um ótimo jeito de
garantir que os endpoints do backend realmente estavam funcionando conforme
esperado.

Outros aprendizados, agora mais técnicos, que consegui retirar desta Ages
foram, um bom conhecimento de FastAPI, a ponto de conseguir criar uma \ac{rest} \ac{api},
implementar uma estrutura de camadas e mexer com o \ac{cors}, um conhecimento
bom de banco de dados, trabalhando em cima de um banco de terceiros, decidindo
quais colunas manter e quais descartar, sempre levando em conta que performance
seria um problema e uma expansão do conhecimento em frontend, podendo
comparar minhas experiências passadas com React \cite{react-docs}, com os atuais códigos feitos
em Angular e tirando as diferenças como aprendizado.

Por fim, também tive uma boa aula de Docker e deploys com um de meus
Ages 3, que provavelmente irá se provar útil quando for minha vez de ser o arquiteto
de software do time, além de uma ótima experiência de interação de grupo, ao
resgatar colegas que se encontravam um tanto quanto perdidos, permitindo que
voltassem a produzir trabalho para o time e aprendessem conhecimentos
apropriados para Ages 1.

Acredito que ter um time tão flexível e produtivo, permitiu que eu adquirisse
conhecimentos básicos de diferentes segmentos que não necessariamente estavam
previstos para um Ages 2 e que certamente irão se provar como um passo adiante
em minhas próximas Ages, bem como no futuro de minha carreira. Ao mesmo tempo
achei muito importante aprender os conceitos de banco de dados e trabalhar com
conexão, integração e build da imagem local do docker, bem como dumps e uso de
um sistema gerenciador de banco de dados, o DBeaver em meu caso.

Saio desta Ages ansioso para cursar matérias como Projeto e Arquitetura de
Software e expandir meus conhecimentos, para poder chegar de Ages 3 como um
ótimo companheiro, tanto em soft skills quanto hard skills. Penso que chegarei com
muito mais maturidade como Ages 3, ainda mais tendo em conta que ela é uma
Ages composta por desafios e tomada de decisões difíceis, que mais do que nunca,
afetam a integridade de grande parte do projeto, bem como todo o time.
